\documentclass[UTF8, 10pt]{ctexart}

\usepackage[margin=0.75in]{geometry}
\usepackage{amsmath, amssymb}
\usepackage{booktabs}
\usepackage{hyperref}
\usepackage{enumitem}
\setlist{nosep,leftmargin=*}
\setlength{\parskip}{0.25em}
\setlength{\parindent}{0pt}

\title{\textbf{课程项目提案}\\
面向通信高效分布式优化的可扩展图划分算法}
\author{小组成员：姓名 1，姓名 2，姓名 3}
\date{CS240：算法设计与分析，2026 年春季学期}

\begin{document}

\maketitle

\section{选题说明}

我们选择参考选题中的 \textbf{Topic 2: Community Detection / Graph
Partitioning}，但将应用场景调整为更现代、也更贴近自身研究方向的
\textbf{通信高效分布式优化}。在分布式优化中，大规模优化问题通常由多个智能体
或 工作节点协同求解，各节点需要通过通信图交换局部信息。因此，图结构会直接
影响通信开销、负载均衡和收敛行为。与常见的社交网络社区发现不同，本项目研究
经典图划分算法如何组织大规模分布式计算。该应用与大规模优化、多智能体系统和
分布式机器学习密切相关。

\section{问题形式化}

设 $G=(V,E,w)$ 是一个无向加权图，其中顶点表示数据块、变量块、智能体或计算
任务，边表示两个顶点之间的通信或依赖关系。给定整数 $k$，目标是将顶点集合
$V$ 划分为 $k$ 个互不相交的子集 $V_1,\ldots,V_k$。

优化目标为
\[
    \min_{V_1,\ldots,V_k}
    \sum_{(u,v)\in E,\; u\in V_i,\; v\in V_j,\; i\neq j} w_{uv},
\]
并满足负载均衡约束
\[
    |V_i| \leq (1+\epsilon)\frac{|V|}{k},
    \quad i=1,\ldots,k.
\]

其中，割边权重表示跨分区通信开销，均衡约束表示不同 工作节点之间的计算负载
需要尽量接近。该问题的输入是一个加权图和目标分区数，输出是一个满足均衡约束
的 $k$ 路顶点划分。

\section{相关工作}

图划分是科学计算、并行计算和大规模图处理中的经典算法问题。谱划分方法利用
图拉普拉斯矩阵的特征向量寻找平衡割；Kernighan--Lin 算法通过局部搜索和顶点
交换降低割边数量；METIS 等多层方法通过图粗化、粗图划分和细化提升可扩展性。
分布式网络优化与本项目密切相关：在一致性优化和分布式梯度方法中，各智能体
根据局部目标函数更新变量并与邻居通信，通信拓扑会显著影响收敛速度和通信开销。

我们计划选择 \textbf{Kernighan--Lin 图划分启发式算法} 作为具体复现的论文 /
方法，谱二分方法和贪心均衡划分将作为基线或变体进行比较。相关论文包括：
\begin{itemize}
    \item B. W. Kernighan and S. Lin, ``An Efficient Heuristic Procedure for
    Partitioning Graphs,'' \emph{Bell System Technical Journal}, 1970.
    DOI: \href{https://doi.org/10.1002/j.1538-7305.1970.tb01770.x}
    {10.1002/j.1538-7305.1970.tb01770.x}.
    \item G. Karypis and V. Kumar, ``A Fast and High Quality Multilevel Scheme
    for Partitioning Irregular Graphs,'' \emph{SIAM Journal on Scientific
    Computing}, 1998. DOI: \href{https://doi.org/10.1137/S1064827595287997}
    {10.1137/S1064827595287997}.
    \item A. Nedic and A. Ozdaglar, ``Distributed Subgradient Methods for
    Multi-Agent Optimization,'' \emph{IEEE Transactions on Automatic Control},
    2009. DOI: \href{https://doi.org/10.1109/TAC.2008.2009515}
    {10.1109/TAC.2008.2009515}.
\end{itemize}

\section{算法与技术计划}

我们计划实现并比较以下几类图划分方法：
\begin{enumerate}
    \item \textbf{贪心均衡划分基线。} 逐个将顶点分配到不同分区，在维持分区
    规模均衡的同时尽量降低局部割边代价。
    \item \textbf{谱二分方法。} 计算图拉普拉斯矩阵，利用 Fiedler 向量将图
    二分，并递归使用二分方法得到 $k$ 个分区。
    \item \textbf{Kernighan--Lin 细化。} 从一个初始划分出发，在满足均衡约束
    的前提下，迭代交换分区之间的顶点以降低割值。
\end{enumerate}

计划使用 Python 实现，主要依赖 NumPy、SciPy 和 NetworkX 等标准科学计算工具。
计划中的实验设置包括网格图、Erdos--Renyi 随机图、随机几何图，并可能加入一个
小规模真实网络。我们计划分析每种方法的时间复杂度，并比较割边大小、负载均衡
比例、运行时间，以及图规模增大时的可扩展性。作为可选的应用层实验，我们计划
模拟一致性平均或分布式梯度下降，测量不同划分下的目标函数误差、一致性误差和
通信成本。

\section{项目范围与预期目标}

项目的基本目标是复现 Kernighan--Lin 启发式算法，实现谱二分和贪心基线，并
比较它们的划分质量与运行时间。这可以直接展示经典图算法如何解决现代分布式
优化中的关键瓶颈：通信高效的工作划分。

如果时间允许，我们会加入以下一个或多个拓展：
\begin{itemize}
    \item 考虑非均匀通信代价的加权图划分；
    \item 与 NetworkX 或 METIS 等已有工具库进行比较；
    \item 加入更完整的分布式优化仿真实验，比较不同网络拓扑和通信预算。
\end{itemize}

\section{团队信息}

\begin{tabular}{@{}lll@{}}
\toprule
成员 & 职责 & 预期贡献 \\
\midrule
姓名 1 & 算法实现 & 谱划分方法与实验 \\
姓名 2 & 理论与写作 & 问题形式化和复杂度分析 \\
姓名 3 & 实验评估 & 分布式优化仿真和图表绘制 \\
\bottomrule
\end{tabular}

\end{document}

